\documentclass{article}
\usepackage{float}
\usepackage{graphicx}
\usepackage{physics}
\usepackage{amssymb}
\usepackage{amsthm}

\theoremstyle{plain}
\newtheorem{theorem}{Theorem}[section]
\newtheorem{lemma}[theorem]{Lemma}
\newtheorem{proposition}[theorem]{Proposition}
\newtheorem{corollary}[theorem]{Corollary}

\theoremstyle{definition}
\newtheorem{definition}{Definition}[section]
\newtheorem{fact}{Fact}[section]

\theoremstyle{remark}
\newtheorem{remark}{Remark}[section]
\DeclareMathOperator{\proj}{proj}
\usepackage{bm}
\usepackage[vmargin=1in, hmargin=1in]{geometry}
\usepackage{mathrsfs}
\usepackage{fontspec}
\usepackage{enumitem}
\usepackage{tikz}
\usetikzlibrary{positioning, quotes, angles, decorations.pathreplacing, calc, arrows.meta}
\renewcommand{\figurename}{Fig.}
\usepackage{chngcntr}
\usepackage{caption}
\let\div\relax
\setlength{\captionmargin}{0pt}
\usepackage{subcaption}
\usepackage{cancel}
\numberwithin{equation}{section}
\newcommand{\N}{\mathbb{N}}
\newcommand{\Z}{\mathbb{Z}}
\newcommand{\R}{\mathbb{R}}

% Note: do NOT define \dd, \tr, \abs, \norm, \var or \div here.
% The physics package already provides them.
\DeclareMathOperator*{\argmin}{arg\,min}
\counterwithin{figure}{section}

\title{The Compound Monitor:\\ Algorithm and Mathematical Machinery}
\author{Ryan Safarha \and Andrey Belokovylenko}
\date{}

\begin{document}
\maketitle

\begin{abstract}
\noindent
A complete account of what the monitor computes, in the order it computes it.
The mathematics is elementary: an AR(1) fit, a cumulative sum, two resampling
procedures and a counting argument. What is not elementary, and what the
document is really about, is \emph{where the alarm threshold comes from}. Every
constant is labelled as a convention, a free choice, a consequence of the
method, or a value set by measurement. Section~\ref{sec:measured} gives the
measured results and Section~\ref{sec:constants} the full constant table.
\end{abstract}

\tableofcontents

%======================================================================
\section{Objects and notation}
%======================================================================

\paragraph{The subject.} One person is a triple $(\mathrm{id}, x, t_0)$ where
$x$ is a real-valued series on a \emph{complete} daily index: every calendar
day between the first and last observation appears as a row, and a day with no
reading is $\mathrm{NaN}$ rather than an absent row. The engine enforces this
at construction. $t_0$ is the day the compound starts. Days before $t_0$ are
the \emph{before period}; days from $t_0$ on are the \emph{monitoring period}.

Completeness of the index matters because the AR(1) model of
Section~\ref{sec:baseline} is a statement about \emph{consecutive calendar
days}. If missing days were dropped rather than marked, two readings a week
apart would be treated as adjacent and the fitted autocorrelation would be
wrong.

\paragraph{The metric.} A metric is a tuple
\[
  (\texttt{name},\ \texttt{unit},\ \texttt{works\_direction},\
   \texttt{harm\_direction},\ \texttt{onset\_days}),
\]
with the two directions in $\{\texttt{up},\texttt{down}\}$. Write
\[
  s_{\mathrm{harm}} = \begin{cases} +1 & \texttt{harm\_direction} = \texttt{up}\\ -1 & \text{otherwise}\end{cases}
  \qquad
  s_{\mathrm{works}} = \begin{cases} +1 & \texttt{works\_direction} = \texttt{up}\\ -1 & \text{otherwise.}\end{cases}
\]
The demo metric is resting heart rate in bpm, with
$s_{\mathrm{works}} = -1$ (a recovery compound should lower it) and
$s_{\mathrm{harm}} = +1$, and an onset of 2 days.

\textbf{Nothing in the engine refers to heart rate.} It sees a series and a
metric. These four fields and the two signs are the only metric-specific
quantities in the entire pipeline, which is why the same code runs on sleep
efficiency without modification.

\subsection{The loader boundary}
\label{sec:loader}

The subject object above is the entire interface between the engine and a
dataset. \texttt{engine.py} refers to no file format, no directory layout and no
particular metric: it sees a series and a metric. Everything dataset-specific
lives in a loader, so adding a dataset means writing a reader, not touching any
of the mathematics that follows.

\paragraph{What a loader must produce.} One value per person per day; a start
date; and \emph{a complete daily index}, with NaN for days that have no reading.

The third requirement is the one that bites, and it fails silently. A loader
that simply omits missing days still runs, and the series still looks
reasonable, but its index is no longer a run of consecutive calendar days.
Equation~\eqref{eq:phi} then pairs readings that are a week apart as though they
were adjacent, $\hat\phi$ is wrong, and nothing anywhere raises an error.

For that reason the adapter layer exposes a single constructor,
\texttt{to\_subject()}, which every loader is expected to end in. It maps
sentinel values to NaN, rejects duplicate dates rather than silently keeping
one, reindexes to the complete daily range, and refuses a before period shorter
than \texttt{MIN\_BEFORE}. A loader that uses it cannot make the mistake above.

\paragraph{What does not transfer.} Running is not the same as being
calibrated. Several constants in Section~\ref{sec:constants} were set by
measuring \emph{this} cohort's characteristics rather than its file format:
\texttt{MIN\_USABLE\_DAYS} from its missingness, and
\texttt{REF\_BASELINE\_DAYS}, \texttt{HORIZON} and \texttt{BLOCK\_DAYS} from
records of about $150$ days with $\phi$ between $0.58$ and $0.92$. The measured
results of Section~\ref{sec:measured} belong to this cohort, this season and
this device.

So a new dataset is loaded with a report of what the engine is about to assume
about it --- series length, coverage, longest gap, $\hat\phi$, $n_{\mathrm{eff}}$,
and the scorable-residual count over one horizon --- and a reminder that checks
6, 7, 10 and 11 must be re-run to obtain its own numbers. The threshold itself
additionally requires at least two subjects, since the default limit of
Section~\ref{sec:reflib} is built from other people; with one subject only the
own-baseline comparator is available, and that is the construction measured at
$14\%$.

\paragraph{A guard rather than a footnote.} Because the failure mode here is a
tool that displays a correct-looking rate belonging to different data, the
measured results carry a fingerprint of the cohort they came from: subject ids,
record spans, value counts and start dates. The display refuses to quote a
false-alarm rate whose fingerprint does not match the data currently loaded, and
says so instead. Portability without that guard would make the tool easier to
misuse, not easier to use.

\paragraph{Three questions.} The monitor answers, for one person:
\begin{enumerate}[nosep]
  \item \textbf{What changed?} Before versus after (Section~\ref{sec:ba}).
  \item \textbf{Is it harming them?} A sustained move away from their own
    normal, flagged automatically (Sections~\ref{sec:baseline}--\ref{sec:threshold}).
  \item \textbf{Is it working for them?} An on/off trial
    (Section~\ref{sec:efficacy}).
\end{enumerate}

%======================================================================
\section{What changed: the before/after comparison}
\label{sec:ba}
%======================================================================

Let $b$ be the non-$\mathrm{NaN}$ values before $t_0$ and $a$ those from $t_0$
on (optionally truncated to a window). With $\bar b, \hat\sigma_b$ the sample
mean and standard deviation ($\mathrm{ddof}=1$) of $b$:
\[
  \text{change} = \bar a - \bar b,
  \qquad
  \text{standardised change} = \frac{\bar a - \bar b}{\hat\sigma_b}.
\]
Dividing by the \emph{before} SD is what makes the number portable: ``$+2.1$''
reads as two of this person's usual swings, independent of units and of how
variable they happen to be.

This quantity is displayed with a fixed caveat, because it is the one
calculation here that cannot support a causal claim:

\begin{quote}
A before/after difference alone can't show the compound caused it; season,
training or illness could. The on/off trial is what attributes change to the
compound.
\end{quote}

Section~\ref{sec:measured} quantifies exactly how badly this caveat is needed:
on data where nobody took anything, a third of monitoring windows show a shift
in the mean.

%======================================================================
\section{Phase I: the baseline model}
\label{sec:baseline}
%======================================================================

The framework is statistical process control. Phase I characterises ``in
control'' from a reference period, here the before period; Phase II monitors
new observations against it.

\subsection{Why a time-series model is needed at all}

Resting heart rate today predicts resting heart rate tomorrow. Every
control-chart method assumes independent observations, so the standard remedy
is to fit a time-series model and monitor its residuals. The relevant
magnitude is the following.

\begin{theorem}[Effective sample size]
\label{thm:neff}
For a stationary AR(1) series with marginal variance $\sigma^2$ and
autocorrelation $\phi$,
\[
  \var(\bar y)
  = \frac{\sigma^2}{n^2}\left(n + 2\sum_{k=1}^{n-1}(n-k)\phi^k\right)
  \;\xrightarrow[n \to \infty]{}\;
  \frac{\sigma^2}{n}\cdot\frac{1+\phi}{1-\phi},
\]
so the series carries the information of
$n_{\mathrm{eff}} = n(1-\phi)/(1+\phi)$ independent observations.
\end{theorem}

\begin{proof}
$\var(\bar y) = n^{-2}\sum_{i,j}\operatorname{Cov}(y_i,y_j)
= n^{-2}\sigma^2\sum_{i,j}\phi^{\abs{i-j}}$. Counting the double sum by
diagonals gives $n$ terms with $\abs{i-j}=0$ and $2(n-k)$ terms for each
$k = 1,\dots,n-1$. Dividing by $n$ and letting $n \to \infty$,
\[
  1 + 2\sum_{k=1}^{n-1}\left(1-\tfrac{k}{n}\right)\phi^k
  \longrightarrow 1 + \frac{2\phi}{1-\phi} = \frac{1+\phi}{1-\phi},
\]
by dominated convergence, since $\sum\phi^k$ converges.
\end{proof}

Measured on this dataset, $\phi$ ranges from $0.58$ to $0.92$ across
participants. At $\phi = 0.85$, $n_{\mathrm{eff}}/n = 0.082$: a 60-day
baseline carries about five independent observations, and treating those 60
days as independent understates standard errors by $\sqrt{12.2} \approx 3.5$.

\begin{remark}[The asymptotic formula fails for short blocks]
\label{rem:short}
The limit in Theorem~\ref{thm:neff} requires $n \gg (1-\phi)^{-1}$. At
$\phi = 0.85$ that is $n \gg 6.7$, so for a block of $m=5$ days the exact sum
must be used. Exactly, $\var(\bar y) = 0.785\,\sigma^2$, giving
$\operatorname{SD}(\bar y) = 0.89\,\sigma$, whereas the asymptotic expression
gives $1.56\,\sigma$ and overstates the noise by a factor of $1.76$. This
matters in Section~\ref{sec:gates}, where blocks are five usable days long.
\end{remark}

\subsection{Fitting}

From the before period, keeping $\mathrm{NaN}$ rows so that adjacency is
preserved, let $P$ be the set of present days.

\paragraph{Location and scale.}
\[
  \hat\mu = \frac{1}{\abs{P}}\sum_{t \in P} x_t,
  \qquad
  \hat\sigma_x^2 = \frac{1}{\abs{P}-1}\sum_{t \in P}(x_t - \hat\mu)^2 .
\]
If $\abs{P} < \texttt{MIN\_BEFORE} = 14$ the engine refuses to fit.

\paragraph{Autocorrelation, with gaps.} Let
$A = \{t : t, t-1 \in P\}$ be the consecutive-present pairs. Then
\begin{equation}
\label{eq:phi}
  \hat\phi = \mathrm{clip}_{[-0.95,\,0.95]}\!
  \left(
  \frac{\displaystyle\frac{1}{\abs{A}}\sum_{t \in A}(x_t-\hat\mu)(x_{t-1}-\hat\mu)}
       {\displaystyle\frac{1}{\abs{P}}\sum_{t\in P}(x_t-\hat\mu)^2}
  \right).
\end{equation}
This is not the textbook lag-1 autocorrelation and the difference is
deliberate: the numerator averages only over pairs that genuinely are one day
apart, while the denominator uses every present day, because every present day
carries information about the variance. On a gapless series the two agree up
to an $O(1/n)$ edge term. The clip keeps the model stationary when a short,
noisy baseline produces $\abs{\hat\phi}$ near or above one.

\paragraph{Residuals, then their scale, then standardisation.} The order
matters and is fixed:
\begin{equation}
\label{eq:resid}
  e_t = (x_t - \hat\mu) - \hat\phi\,(x_{t-1}-\hat\mu),
  \qquad
  \hat\sigma_e^2 = \frac{1}{\abs{E}-1}\sum_{t \in E} e_t^2,
  \qquad
  z_t = \frac{e_t}{\hat\sigma_e},
\end{equation}
where $E = A$ is the set of $t$ for which $e_t$ is defined. If either $x_t$ or
$x_{t-1}$ is missing then $e_t$ is $\mathrm{NaN}$; the first day of the series
has no predecessor and is $\mathrm{NaN}$ too.

Under the fitted model the $z_t$ are approximately independent with mean $0$
and variance $1$, which is what the rest of the pipeline assumes. Whether that
survives contact with real post-start data is measured, not asserted: see
Section~\ref{sec:diagnostic}.

\paragraph{Carrying the fit forward.} The monitoring period is standardised
with the \emph{same} $\hat\mu, \hat\phi, \hat\sigma_e$ fitted on the before
period, applied along the whole series, so the first monitored day uses the
last before-day as its predecessor and no day is wasted at the join.

\paragraph{Valid steps.} For any window, write
\[
  n_{\mathrm{valid}} = \#\{t \in \text{window} : z_t \ne \mathrm{NaN}\}.
\]
Because $e_t$ needs two consecutive present days, if days are missing roughly
at random with rate $q$ then $n_{\mathrm{valid}} \approx (1-q)^2$ times the
calendar length. At $q = 0.34$ a 30-day window supplies about 13 valid steps,
not 30. This number drives the calibration in Section~\ref{sec:threshold} and
is the single most frequent source of silent error in a design like this.

%======================================================================
\section{The detector}
\label{sec:detector}
%======================================================================

\begin{definition}[One-sided CUSUM]
\label{def:cusum}
Given standardised residuals $z_t$, a reference value $k>0$ and a sign
$s = s_{\mathrm{harm}}$,
\[
  S_0 = 0,
  \qquad
  S_t = \max\!\left(0,\; S_{t-1} + s\,z_t - k\right),
  \qquad \text{alarm when } S_t > h .
\]
If $z_t$ is $\mathrm{NaN}$ then $S_t = S_{t-1}$ and no step is taken.
\end{definition}

$S_t$ is a random walk with a reflecting barrier at the origin. The barrier is
what makes it a monitor rather than an accumulator: evidence that does not
persist is forgotten.

\begin{proposition}[The hypotheses are separated by the sign of a drift]
\label{prop:drift}
If $\mathbb{E}[s z_t] < k$ the increments have negative mean, the walk is
absorbed at the barrier and returns to $0$ infinitely often. If
$\mathbb{E}[s z_t] > k$ the increments have positive mean, and by the strong
law $S_t / t \to \mathbb{E}[s z]-k > 0$, so $S_t$ grows linearly and crosses
any fixed $h$ in expected time $O(h)$.
\end{proposition}

So $k$ is not a threshold but a \emph{tilt}: it is chosen so that in control
the walk drifts down and out of control it drifts up, and detection reduces to
detecting the sign of a drift. This is exactly why the detector sees what a
per-day limit and a self-report both miss. A long run of mildly elevated days
climbs to the same place as one dramatic day, and nobody files a report for
feeling slightly worse for three weeks.

\subsection{Choosing $k$}

The standard rule is to set $k$ at half the shift one wants to detect,
measured \emph{in the units being monitored}. The units being monitored are
residual standard deviations, while the target shift $\Delta$ is expressed in
raw baseline SDs, so the two are related by the following attenuation.

\begin{proposition}[Whitening attenuates a step shift]
\label{prop:atten}
A sustained shift of $\delta\sigma_x$ in the mean of an AR(1) series appears in
the standardised residuals as
\[
  \delta\sqrt{\frac{1-\phi}{1+\phi}} .
\]
\end{proposition}

\begin{proof}
Once the shift has persisted for more than one day, both $x_t$ and $x_{t-1}$
are displaced by $\delta\sigma_x$, so by \eqref{eq:resid} the residual retains
$(1-\phi)\delta\sigma_x$. Dividing by
$\sigma_e = \sigma_x\sqrt{1-\phi^2}$,
\[
  \frac{(1-\phi)\delta\sigma_x}{\sigma_x\sqrt{(1-\phi)(1+\phi)}}
  = \delta\sqrt{\frac{1-\phi}{1+\phi}} . \qedhere
\]
\end{proof}

Hence
\begin{equation}
\label{eq:k}
  k = \frac{\Delta}{2}\sqrt{\frac{1-\hat\phi}{1+\hat\phi}},
  \qquad \Delta = 2.0,
\end{equation}
computed per person from that person's own $\hat\phi$. At $\hat\phi = 0.85$
the attenuation factor is $0.28$, so $k \approx 0.28$ and a $2\sigma$ raw
shift arrives as about $0.56$ in residual units.

\begin{remark}[The price, and a prediction that was wrong]
Proposition~\ref{prop:atten} says the detector sees only $28\%$ of the shift it
would see on independent data. That is the accepted price of a valid
false-alarm rate, and an invalid false-alarm rate is not a weaker guarantee but
no guarantee at all. It is worth recording that the author predicted whitening
would \emph{cost power} relative to an unwhitened chart, and it does not: the
null is attenuated by the same factor, so the signal-to-noise ratio is
unchanged. Only the ratio matters.
\end{remark}

%======================================================================
\section{The threshold}
\label{sec:threshold}
%======================================================================

Let
\[
  M = \max_{1 \le t \le n_{\mathrm{valid}}} S_t
\]
be the running maximum over one monitoring course. We need
\[
  h = F_M^{-1}(1-\alpha), \qquad \alpha = \texttt{FP} = 0.05,
\]
the $(1-\alpha)$ quantile of the null distribution of $M$. There is no closed
form. \textbf{Where $F_M$ comes from is the entire design question of this
project}, and the false-alarm rate is the only quantity the project really
measures.

\begin{proposition}[The horizon must match]
\label{prop:horizon}
$F_M$ depends on the number of \emph{scored steps}, and $M$ is nondecreasing in
that number, since the maximum is taken over a larger set. A limit calibrated
for $n_1$ steps and applied over $n_2 > n_1$ over-alarms; applied over
$n_2 < n_1$ it under-alarms.
\end{proposition}

This is why $n_{\mathrm{valid}}$, not the calendar length, is passed to every
calibration routine. A person with 34\% missing days takes about 13 steps in a
30-day window, and a threshold built for 30 steps would make them look
well-calibrated while silently costing sensitivity --- on exactly the subjects
whose data is worst.

\subsection{Calibration A: the own-baseline bootstrap (the comparator)}

\begin{enumerate}[nosep]
  \item Take the standardised before-period residuals $z^{\mathrm{base}}$,
    $\mathrm{NaN}$s dropped.
  \item Draw $n_{\mathrm{valid}}$ of them with replacement; run
    Definition~\ref{def:cusum}; record $\max_t S_t$.
  \item Repeat $N = 1000$ times. Set $h$ to the $95$th percentile of those
    maxima.
\end{enumerate}

Because Section~\ref{sec:baseline} has already whitened the series, resampling
single residuals with replacement is the ordinary residual bootstrap; no block
structure is needed here.

This calibration needs nobody but the subject, which is its only advantage.

\begin{proposition}[Why it fails]
\label{prop:fails}
The procedure computes a quantile of $f^{\mathrm{assumed}}_{H_0}$, the null
built from the before period. The monitoring days follow
$f^{\mathrm{true}}_{H_0}$. If the mean shifts between the two periods these are
different distributions, and the $(1-\alpha)$ quantile of one is not that of
the other. Increasing $N$ converges to the quantile of the wrong distribution.
\end{proposition}

The error is therefore bias, not noise, and it has a direction: the assumed
null contains none of the post-baseline drift, so it is too tight, $h$ is too
low, and false alarms exceed nominal. Measured: $14.0\%$ against a nominal
$5\%$.

\subsection{Calibration B: the reference library (the default)}
\label{sec:reflib}

Replace the assumed null with a real one, built from other people who took
nothing.

\paragraph{Building the library.} For each subject $j$ other than the target,
slide through $j$'s post-start series \emph{in valid-step space}: for every
index $i$ whose residual $z_i$ is defined and which lies at or after $j$'s own
start,
\begin{enumerate}[nosep]
  \item \textbf{Rebase.} Form the subject restricted to the
    $\texttt{REF\_BASELINE\_DAYS} = 60$ calendar days ending at $i-1$, with
    start $i$.
  \item \textbf{Re-fit.} Fit $\hat\mu,\hat\phi,\hat\sigma_e$ on that
    60-day window. If it holds fewer than $14$ present days, skip this window.
  \item \textbf{Score.} Compute residuals from day $i-1$ onward, standardise,
    drop $\mathrm{NaN}$s and keep the first $\texttt{HORIZON}=30$ of them.
\end{enumerate}
One window is produced per valid post-start day. On this dataset that is $876$
windows from $13$ subjects.

\paragraph{The threshold.} Given the target's $n_{\mathrm{valid}}$ and $k$:
take every library window belonging to another subject that supplies at least
$n_{\mathrm{valid}}$ residuals, truncate each to its first $n_{\mathrm{valid}}$,
run the CUSUM \emph{with the target's own $k$}, and set $h$ to the $95$th
percentile of the resulting maxima. The statistic is therefore ``what would
\emph{this} detector have done on \emph{that} person's unexposed data,'' which
is the right null for the question being asked.

Leaving the target out is not optional. Including their own data would fit the
decision rule to the test data and the error estimate would mean nothing.

\paragraph{Why the baseline is re-fitted per window.} A single fixed baseline
would score a window 90 days later against a 90-day-old reference, which no
real course does. Worse, most subjects in this dataset share a start date, so
the earliest windows all cover the same calendar stretch; sliding one step at a
time then replicates that one shared drift episode into hundreds of
overlapping windows until the $95$th percentile lands inside it. Measured: the
demo subject's $h$ rises from $8.5$ to $16.9$ and detection falls to $21\%$.
Re-fitting per window removes it.

\subsection{What the library can and cannot buy}
\label{sec:quantile}

\begin{proposition}[A quantile from a small sample]
\label{prop:smallq}
The sample $(1-\alpha)$ quantile from $m$ independent draws from a density $f$
has asymptotic variance $\alpha(1-\alpha)/(m f(q)^2)$, which is large for small
$m$ and larger still in a tail where $f(q)$ is small. At $\alpha = 0.05$ and
$m = 13$, $\Pr(\text{all 13 draws lie below } q) = 0.95^{13} = 0.51$: about
half the time the entire sample sits below the quantity being estimated.
\end{proposition}

Overlapping windows do not increase $m$. Windows from one person share almost
all their data, so the effective sample size is the \emph{subject} count, not
the window count. Both are reported everywhere.

\begin{remark}[Measured]
Recomputing $h$ with each library subject left out in turn moves it by
$10\%$ to $46\%$ of its value (median $37\%$); for the demo subject the range
is $5.75$ to $8.68$. One subject dominates: $22$ of the $30$ largest library
maxima are theirs, and dropping them lowers $h$ for $10$ of $13$ subjects.

Yet the false-alarm \emph{rate} barely moves: $7.0\%$ for twelve of the
thirteen removals, $11.6\%$ for the one. The threshold swings across a gap in
the null's order statistics where few observed maxima lie, so a large change in
$h$ produces a small change in the rate. Both numbers are reported, because the
first is the measured reason a deployment needs a reference cohort of $50$ to
$100$ people rather than $13$, and enlarging the library cannot substitute:
going from $60$ to $500$ windows moved $h$ by $0.4$.
\end{remark}

%======================================================================
\section{Is it working? The on/off trial}
\label{sec:efficacy}
%======================================================================

The safety arm needs no assumption about the compound. The efficacy arm needs
to attribute a change to it, which requires a design rather than an analysis.

\paragraph{The design.} $n$ cycles. Each cycle is two adjacent blocks of
$\texttt{BLOCK\_DAYS}=7$ days, one on and one off, with the order inside the
cycle decided by a coin flip. The first $\texttt{onset\_days}$ of every block
are excluded from analysis as washout, since a compound with a half-life is
still present after a block ends. With a 2-day onset, each block contributes
five usable days.

\paragraph{The statistic.} For each cycle with both blocks usable (at least
$\texttt{MIN\_USABLE\_DAYS}=2$ post-washout days each), let
\[
  d_c = \overline{(\text{on block})} - \overline{(\text{off block})},
\]
the difference of post-washout means. Block means rather than individual days,
because days are autocorrelated and treating them as independent inflates false
positives; averaging within blocks is the standard n-of-1 remedy.

\subsection{The randomisation test}

\begin{definition}[Randomisation $p$-value]
Let $G$ be the set of label assignments the randomisation could have produced
and $T(y,g)$ the statistic under assignment $g$. Then
\[
  p = \frac{\#\{g \in G : \abs{T(y,g)} \ge \abs{T(y,g_{\mathrm{obs}})}\}}{\abs{G}} .
\]
\end{definition}

Validity comes from the coin flips that were actually made. Under the null that
the compound does nothing, the observed data would have been identical
whichever assignment came up, so all $\abs{G}$ values of $T$ were equally
likely and the rank of the observed one is uniform. \textbf{No normality, no
independence and no stationarity is assumed, so none of them can break it.}
This is why the efficacy arm is untouched by the $\phi \approx 0.85$ problem
that dominates the safety arm.

Here the order is randomised \emph{within each cycle}, so $G$ is the set of
sign patterns $\{\pm 1\}^n$ and $\abs{G} = 2^n$.

\begin{theorem}[Minimum attainable $p$]
\label{thm:pmin}
The map $g \mapsto g^c$ that flips every sign is an involution on $G$ with
$T(y,g^c) = -T(y,g)$. Hence $\abs{T}$ is constant on complementary pairs, the
observed assignment always ties with its own complement, the numerator is at
least $2$, and
\[
  p_{\min} = \frac{2}{2^n}.
\]
\end{theorem}

\begin{corollary}[$\texttt{MIN\_CYCLES} = 6$]
$p_{\min} < \alpha = 0.05$ requires $2^n > 40$, hence $n \ge 6$. At $n=3$,
$p_{\min} = 0.25$; at $n=5$, $0.0625$; at $n=6$, $0.03125$.
\end{corollary}

This is a statement about the cardinality of a finite set, not about power. No
effect size, no amount of data and no better instrument changes it. Below six
cycles the app refuses the schedule and says why, \emph{before} it is run,
because a tool that accepts a 3-cycle schedule, runs it for six weeks and
reports $p = 0.25$ has produced a confident-looking null determined by the
calendar rather than the compound.

\begin{corollary}[The test is conservative by construction]
\label{cor:conservative}
The attainable $p$-values lie on the grid $\{2/2^n, 4/2^n, \dots\}$. At $n=6$,
$p < 0.05$ means $p = 2/64 = 0.031$ exactly, so the true false-positive rate is
about $3.1\%$ rather than $5\%$. Measured: $3.6\%$ over $140$ null runs.
\end{corollary}

\subsection{Dropped cycles}

If some cycles are unusable the test runs on the survivors, and the
implementation handles this by zero-padding rather than by re-indexing.

\begin{lemma}[Zero-padding preserves the $p$-value]
\label{lem:pad}
Let $k$ of $n$ cycles be usable and set $d_c = 0$ for the rest. Then the
proportion of sign patterns in $\{\pm1\}^n$ at least as extreme as the observed
one equals the proportion over the $2^k$ patterns on the usable cycles alone.
\end{lemma}

\begin{proof}
The statistic $\sum_c \varepsilon_c d_c$ does not depend on $\varepsilon_c$
when $d_c = 0$. So each of the $2^k$ distinct values arising from the usable
cycles is attained by exactly $2^{n-k}$ patterns of the full set. Both
numerator and denominator are multiplied by $2^{n-k}$, and the ratio is
unchanged.
\end{proof}

The reported floor uses $k$, not $n$: a trial that was scheduled for eight
cycles and kept seven is judged at $p_{\min} = 2/2^7$.

\subsection{The verdict}

With $k$ usable cycles, $\bar d$ their mean and $p$ the randomisation
$p$-value:
\[
  \begin{cases}
    \textbf{Not enough usable cycles} & k < \texttt{MIN\_CYCLES} \\
    \textbf{Consistent effect: working} & p < \alpha,\ \operatorname{sign}\bar d = s_{\mathrm{works}} \\
    \textbf{Consistent effect: opposite to intended} & p < \alpha,\ \operatorname{sign}\bar d = -s_{\mathrm{works}} \\
    \textbf{No consistent effect detected} & \text{otherwise.}
  \end{cases}
\]
The last is always shown with the sentence \emph{``This is not the same as no
effect. This schedule reliably detects effects of about $m$ SDs or larger,''}
with $m$ from Section~\ref{sec:gates}. An ``opposite'' verdict is displayed
alongside the harm status, since a compound moving the metric the wrong way is
a safety observation as much as an efficacy one.

%======================================================================
\section{The three gates}
\label{sec:gates}
%======================================================================

Before a schedule is run, three numbers are computed and shown. They answer
three different questions, and conflating them is how a trial ends up producing
a meaningless null.

\begin{center}
\begin{tabular}{lll}
quantity & question & section \\
\hline
$p_{\mathrm{floor}}$ & can a result exist at all? & \ref{thm:pmin} \\
$p_{\mathrm{complete}}$ & will the trial finish? & \ref{sec:pcomplete} \\
$m_{\mathrm{det}}$ & would a real effect be found? & \ref{sec:mdet} \\
\end{tabular}
\end{center}

\subsection{A precondition before any probability}
\label{sec:fit}

If the schedule runs past the last day with data, there is nothing to compute.
A schedule of $n$ cycles beginning at $t_0$ ends at
$t_0 + 2n\cdot\texttt{BLOCK\_DAYS} - 1$, and the engine raises an error naming
how many cycles would fit instead. This is a length check of the same kind the
loader already applies to the before period, and it is the right answer for a
device that stopped recording: not ``$81\%$ chance of completing'' but ``this
schedule does not fit in the available record.''

On this dataset, six of thirteen subjects are rejected here and seven fit.

\subsection{$p_{\mathrm{complete}}$: will the trial finish?}
\label{sec:pcomplete}

A block is dropped if it has fewer than $\texttt{MIN\_USABLE\_DAYS}$
post-washout days with data; a cycle needs both of its blocks; the verdict needs
$\texttt{MIN\_CYCLES}$ cycles. So
\[
  p_{\mathrm{complete}} = \Pr\big(\text{at least } \texttt{MIN\_CYCLES}
  \text{ of } n \text{ cycles survive}\big).
\]

\paragraph{Not the probability that all survive.} The verdict needs six usable
cycles, not every scheduled one: schedule ten, lose three, and the test runs on
seven. The all-must-survive version \emph{decreases} as cycles are added, which
would make ``consider more cycles'' actively harmful advice; the correct version
increases. The two coincide exactly at $n = \texttt{MIN\_CYCLES}$, which is why
the error is invisible there.

\paragraph{The binomial version, and why it is not enough.} If missing days
were independent with coverage $c$, then with $u$ post-washout days per block,
\[
  p_{\mathrm{block}} = \Pr(\mathrm{Bin}(u,c) \ge \texttt{MIN\_USABLE\_DAYS}),
  \qquad
  p_{\mathrm{complete}} = \Pr\!\left(\mathrm{Bin}(n, p_{\mathrm{block}}^2) \ge \texttt{MIN\_CYCLES}\right).
\]
Real missing days come in runs, and a run is far more destructive than the same
number of scattered days: losing two days of five leaves the block usable,
while losing five consecutive days destroys the block, the cycle and (without
slack) the trial. The binomial cannot see this. Measured, the binomial
predicted $0.06$ to $0.90$ for the two gap-prone subjects, and both reached a
verdict in $0$ of $20$ real schedules.

\paragraph{The moving-block bootstrap.} Instead of modelling the gap pattern,
resample it. Take the subject's own before-period missingness indicator
$m \in \{0,1\}^{N}$. For block length $L$, draw $\lceil T/L \rceil$ starting
positions uniformly from $\{0,\dots,N-L\}$, concatenate the blocks $m[s:s+L]$
and truncate to the trial length $T$. Lay the result over the planned schedule
and count whether at least $\texttt{MIN\_CYCLES}$ cycles survive. Repeat
$10{,}000$ times.

Runs appear automatically because runs are what is being resampled, and the
pattern is the subject's own rather than an assumed one.

\paragraph{The block length is not a free choice.} $L$ changes the answer: for
one subject $p_{\mathrm{complete}}$ is $0.90 / 0.79 / 0.64 / 0.77$ at
$L = 3,5,7,14$, and the warning fires at $0.8$. A gate whose output depends on
an arbitrary constant is not a gate, so the engine computes all four and
\textbf{reports the minimum}. That is the conservative end of a sensitivity
analysis rather than a tuned value: a warning should fire if any defensible
assumption says it should. All four are shown in the app.

The non-monotonicity in $L$ is itself informative. It means that subject's gaps
have structure at a particular scale rather than being random runs, and 7-day
blocks happen to reproduce it.

If the before period is shorter than $\texttt{MIN\_RESAMPLE\_DAYS} = 28$ there
is too little pattern to resample and the binomial is used, flagged as
optimistic.

\paragraph{One standing caveat.} $p_{\mathrm{complete}}$ is conditional on the
device continuing to record as it has been. Nothing can predict a watch that
stops.

\subsection{$m_{\mathrm{det}}$: would a real effect be found?}
\label{sec:mdet}

$p_{\mathrm{floor}}$ says a result is \emph{possible}; this says whether it is
\emph{likely}. Simulate the whole pipeline:

\begin{enumerate}[nosep]
  \item Generate $N$ trials of length $T = 2nB$ from the fitted AR(1) with unit
    marginal variance: $x_1 \sim N(0,1)$ and
    $x_t = \hat\phi x_{t-1} + \varepsilon_t$,
    $\varepsilon_t \sim N(0, 1-\hat\phi^2)$. Unit marginal SD means an effect
    of $\delta$ is $\delta$ baseline SDs by construction.
  \item Apply a missingness draw from Section~\ref{sec:pcomplete}, at the block
    length that produced the minimum.
  \item Randomise the on/off order within each cycle.
  \item Add $s_{\mathrm{works}}\,\delta$ on the on-blocks, ramped linearly over
    $\texttt{onset\_days}$ at the start of each on-block --- the same shape
    \texttt{plant\_effect} uses, so ``would be found'' and ``was found'' cannot
    contradict each other.
  \item Compute block means, cycle differences, the usable-cycle mask and the
    exact randomisation $p$. A trial counts as detected when it has at least
    $\texttt{MIN\_CYCLES}$ usable cycles, $p < \alpha$, \emph{and} the mean
    difference has the intended sign.
  \item Sweep $\delta$ over $0.1, 0.2, \dots$ and report the first reaching
    $80\%$.
\end{enumerate}

\paragraph{The plateau rule.} Power can never exceed $p_{\mathrm{complete}}$,
since a trial that does not finish cannot detect anything. When that ceiling
sits within two standard errors of the target,
\[
  p_{\mathrm{complete}} < 0.80 + 2\sqrt{\frac{0.8 \times 0.2}{N}} \approx 0.825
  \quad (N = 1000),
\]
the power curve is flat near the target and any crossing is a draw from a flat
region. The engine then reports \textbf{not reachable} rather than a number,
because quoting $2.4$ when $2.8$ and $3.2$ are equally consistent implies a
precision that does not exist.

%======================================================================
\section{Planting}
%======================================================================

PMData participants took nothing, so there is no real effect in the data. Two
functions inject one, and they are the only way an effect enters.

\paragraph{Harm.} From $t_0 + 5$ days, add
$s_{\mathrm{harm}}\,\delta\,\hat\sigma_x$ to every value, ramped linearly over
three days and sustained thereafter. $\mathrm{NaN}$ days stay $\mathrm{NaN}$.

\paragraph{Effect.} Add $s_{\mathrm{works}}\,\delta\,\hat\sigma_x$ during
on-blocks only, ramped over $\texttt{onset\_days}$ at the start of each
on-block, matching the washout exactly so that the ramp falls entirely inside
the excluded days.

Both return a new subject with $\texttt{simulated} = \mathrm{True}$, a flag
that survives every transformation and drives the on-screen label. The two
controls are mutually exclusive in the app, since the works and harm directions
are opposite and planting both would partially cancel during on-blocks.

\begin{remark}[Planting is a measurement, not a fabrication]
Injecting a known signal to characterise an instrument is what calibration is:
a smoke detector is tested with a known quantity of smoke. The property that
makes it honest is that the answer is known before you look, so there is no
room to interpret the outcome favourably. Note also the separation of concerns:
the \emph{false-alarm} rate is measured on entirely unmodified real data, where
every alarm is false by construction, and planting enters only the
\emph{detection} measurement. Finally, the result is unflattering --- $53\%$ ---
which is what a measurement looks like and what a fabrication does not.
\end{remark}

%======================================================================
\section{The acceptance checks}
%======================================================================

Eleven checks, run as one script. Some have pass conditions; three are
measurements reported without a threshold, because a quantity you are
measuring cannot also be a target you are hitting.

\begin{center}
\begin{tabular}{rll}
\# & what it tests & kind \\
\hline
1 & loader: complete index, no zeros, NaN accounting & pass/fail \\
2 & before/after on a known $+5$ shift at SD $2$ & pass/fail \\
3 & $\hat\phi$ recovery and residual whitening on synthetic AR(1) & pass/fail \\
4 & CUSUM at rest: all-zero $z$ gives all-zero $S$ & pass/fail \\
5 & bootstrap self-consistency: fresh windows alarm at $5\%$ & pass/fail \\
6 & \textbf{false alarms on real unexposed data, both limits} & measured \\
7 & detection with a planted $2\sigma$ harm, both limits & pass/fail \\
8 & the three gates, the length precondition, block sensitivity & pass/fail \\
9 & efficacy wiring on a planted effect & pass/fail \\
10 & \textbf{efficacy false positives on real data} & measured \\
11 & \textbf{baseline-model diagnostic on real monitoring days} & measured \\
\end{tabular}
\end{center}

Check 5 deserves a note: it resamples from the same residuals the limit was
built from, so it verifies the \emph{code} and says nothing about real-world
validity. It is checks 6 and 11 that do that.

\subsection{Check 11: the free diagnostic}
\label{sec:diagnostic}

Everything downstream assumes the fitted $\hat\mu, \hat\phi, \hat\sigma_e$
still describe the person during monitoring. Check 11 tests that directly, with
no detector, no simulation and no planting. On each real monitoring window,
compute $z_t$ and report the distribution across windows of
\[
  \bar z \quad (\text{should be } 0),
  \qquad
  \operatorname{SD}(z) \quad (\text{should be } 1),
\]
together with the share of windows where
$\abs{\bar z} > 2/\sqrt{n_{\mathrm{valid}}}$, roughly two standard errors.

This is the free diagnostic: no labels, no held-out data, no detector output.
It separates the two ways the model can fail, which nothing else in the build
does. $\bar z \ne 0$ is a \emph{location} error, the mean moved;
$\operatorname{SD}(z) \ne 1$ is a \emph{scale} error, $\hat\phi$ or
$\hat\sigma_e$ was mis-estimated. The $14\%$-versus-$7\%$ comparison tells you
the null is wrong but not which.

%======================================================================
\section{Measured results}
\label{sec:measured}
%======================================================================

PMData (Thambawita et al., Simula; CC BY 4.0), daily resting heart rate from 16
participants of whom 13 are usable, recorded November to March. Nobody took
anything, so every alarm below is false by construction. Windows start at days
$60, 74, 88, 102$, each with its own 60-day baseline; $43$ windows from $13$
subjects.

\paragraph{Safety arm.}
\begin{center}
\begin{tabular}{lcc}
 & reference limit (default) & own baseline (comparator) \\
\hline
false alarms & $7.0\%$ ($3/43$) & $14.0\%$ ($6/43$) \\
Wilson 95\% CI & $2.4$--$18.6\%$ & $6.6$--$27.3\%$ \\
subjects with an alarm & $2/13$ & $5/13$ \\
$2\sigma$ harm caught within 30 d & $53\%$ ($23/43$) & $79\%$ ($34/43$) \\
days to alarm, median (IQR) & $13$ ($10$--$22$) & $11$ ($8$--$17$) \\
\end{tabular}
\end{center}

\noindent Intervals treat windows as independent; they overlap within a
person, so the true intervals are wider.

The own-baseline limit catches more \emph{because its $h$ is lower}, which is
the same fact as its $14\%$ false-alarm rate rather than a separate benefit.
These are one detector at two points on the same curve, not two methods.

\paragraph{Detection depends on the person's own drift.} Splitting the 43
windows by their check-11 mean $z$:
\begin{center}
\begin{tabular}{lcc}
window's own drift & planted harm caught & median day \\
\hline
toward harm ($\bar z > 0$) & $14/15$ ($93\%$) & $11$ \\
away from harm & $9/28$ ($32\%$) & $21$ \\
\end{tabular}
\end{center}
So $53\%$ is not a per-person coin flip; it is a blend of near-certain and
unlikely, and which applies cannot be known in advance. The split uses the same
days the detector scores, so it describes outcomes rather than predicting them.
It also cannot be turned into a runtime confidence qualifier: a window drifting
toward harm and a window containing harm are the same observation, so $\bar z$
carries no information the detector does not already use.

\paragraph{Efficacy arm.} False positives $3.6\%$ ($5/140$ runs, Wilson
$1.5$--$8.1\%$), against $3.1\%$ expected from
Corollary~\ref{cor:conservative}. Minimum detectable effect $1.7$ to
$1.9$ baseline SDs for subjects whose schedule fits, about $\pm 0.1$ of Monte
Carlo noise. Seven of thirteen subjects can run a six-cycle trial; none fits
eight.

\paragraph{Check 11.}
\begin{center}
\begin{tabular}{lcc}
 & measured over 43 windows & if the model held \\
\hline
$\operatorname{SD}(z)$ & median $0.97$ (IQR $0.79$--$1.10$), RMS $1.02$ & $1$ \\
$\bar z$ & median $-0.09$ (IQR $-0.43$ to $+0.10$), range $-0.85$ to $+0.69$ & $0$ \\
$\abs{\bar z} > 2/\sqrt{n_{\mathrm{valid}}}$ & $15/43$ ($35\%$), Wilson $22$--$50\%$ & about $5\%$ \\
lag-1 autocorr.\ of $z$ & median $+0.12$ (IQR $-0.08$ to $+0.27$) & $0$ \\
\end{tabular}
\end{center}

Read this as three statements.

\textbf{The scale holds.} $\hat\phi$ and $\hat\sigma_e$, fitted on 60 days
before exposure, carry over to the monitoring period. This is the half of the
model that had to work: a scale error would mean $z$ is not comparable between
people, so borrowing other people's CUSUM maxima would be borrowing numbers in
different units and Section~\ref{sec:reflib} would be invalid.

\textbf{The location fails, with no common direction.} The median $\bar z$ is
essentially zero while the spread runs from $-0.85$ to $+0.69$. This is the
measured answer to ``why not just correct for the drift?'': there is no common
drift to subtract, so a population template would change nothing. Only a
reference-group threshold helps, because it absorbs spread where a template
removes a mean.

\textbf{A small residual autocorrelation remains.} The $2/\sqrt{n}$ bound
assumes independent residuals, and $z$ retains a lag-1 of $+0.12$, consistent
with $\hat\phi$ being biased low by roughly $(1+3\phi)/n \approx 0.06$ from a
60-day fit. Widening the bound to allow for it still leaves $13$ of $43$
windows. A bias-corrected $\hat\phi$ is the measured next step; it is not
implemented, and the screens state the measured $7\%$ rather than the nominal
$5\%$, so nothing is claimed that has not been measured.

%======================================================================
\section{The constants}
\label{sec:constants}
%======================================================================

\begin{center}
\begin{tabular}{llp{7.6cm}}
name & value & status \\
\hline
$\texttt{FP}$ & $0.05$ & convention; false-alarm rate per course \\
$\texttt{ALPHA}$ & $0.05$ & convention; efficacy significance level \\
$\texttt{N\_SIMS}$ & $1000$ & enough for a stable 95th percentile \\
$\texttt{HORIZON}$ & $30$ d & choice; length of one course \\
$\texttt{DELTA}$ & $2.0$ & choice; shift the detector is tuned for \\
$\texttt{MIN\_BEFORE}$ & $14$ d & choice; refuse to fit below this \\
$\texttt{BLOCK\_DAYS}$ & $7$ & choice; valid only for compounds acting within \texttt{onset\_days} \\
$\texttt{REF\_BASELINE\_DAYS}$ & $60$ d & choice; re-fitted per library window \\
$\texttt{MISS\_BLOCK\_LENGTHS}$ & $3,5,7,14$ & choice; $p_{\mathrm{complete}}$ is the \emph{minimum} over them \\
$\texttt{MIN\_RESAMPLE\_DAYS}$ & $28$ & choice; below this, fall back to the binomial \\
start day & $60$ & choice (demo only; there is no real exposure) \\
plant sizes & $2.0,\ 1.0$ & choice; always labelled \\
\hline
$k$ & eq.~\eqref{eq:k} & \textbf{follows} from $\Delta$, $\hat\phi$ and Prop.~\ref{prop:atten} \\
$\texttt{MIN\_CYCLES}$ & $6$ & \textbf{follows} from $p_{\min}=2/2^n$ and \texttt{ALPHA} \\
$\texttt{DEFAULT\_CYCLES}$ & $8$ & \textbf{follows}: two cycles of slack, so one bad week does not end a trial \\
plateau cutoff & $0.825$ & \textbf{follows}: two standard errors above $0.80$ at $N=1000$ \\
\hline
$\texttt{MIN\_USABLE\_DAYS}$ & $2$ & \textbf{measured}: at $3$, completion is $0.49/0.24/0.05$ at $20/26/34\%$ missing, and dropping to $2$ costs nothing on complete data \\
$\texttt{N\_SIMS\_COMPLETE}$ & $10{,}000$ & \textbf{measured}: $\pm 0.01$ noise, so the $0.8$ warning is not seed-dependent \\
\end{tabular}
\end{center}

%======================================================================
\section{What this does and does not establish}
%======================================================================

\paragraph{Established.} Calibrating a per-person monitor from the subject's
own pre-exposure baseline alarms on $14\%$ of unexposed monitoring courses
against a nominal $5\%$, and the cause is a location error with no common
direction (Section~\ref{sec:diagnostic}). A leave-one-subject-out reference
threshold brings this to $7\%$ while still catching $53\%$ of a sustained
$2\sigma$ shift, typically by day 13. Both numbers are measured on real data
from people who took nothing.

\paragraph{Not established, and stated as such.}
\begin{itemize}[nosep]
  \item Transfer. The $7\%$ and $53\%$ belong to this cohort, this season and
    this device. Check 11 shows the location error is substantial and
    person-specific, so a new deployment must measure its own.
  \item Seasonality. Every participant was recorded November to March, so
    cohort-wide seasonal drift sits inside the empirical null instead of being
    caught by it. The $7\%$ is optimistic for a cohort spread across different
    months, and this cannot be fixed with this dataset.
  \item Cohort size. With $13$ subjects the threshold is effectively set by one
    person's drift episode (Section~\ref{sec:quantile}). A deployment needs
    $50$ to $100$.
  \item Trial length. The efficacy arm needs six cycles at minimum and eight by
    default, roughly $112$ days plus a baseline. No subject here fits eight.
  \item Multiplicity. One primary metric per compound, fixed in advance.
    Monitoring four at a nominal $5\%$ each gives $1-0.95^4 = 18.5\%$ per
    course.
  \item Portability. The engine runs on any dataset that can be turned into
    subjects (Section~\ref{sec:loader}), and this has been exercised once
    already on a second metric with different units, a different direction
    convention and a different missingness pattern. But the constants and the
    measured numbers are this cohort's, and the display refuses to quote a rate
    it has not measured on the data in front of it.
\end{itemize}

\paragraph{Hypotheses tested and rejected.} That AR(1) whitening costs power
(it does not; the null attenuates identically). That a larger reference library
lowers $h$ (twelve times more windows moved it by $0.4$). That drift could be
removed with a population template on this dataset (there is no common drift to
remove). Each was predicted, measured and abandoned.

\begin{thebibliography}{9}
\bibitem{alwan} L. C. Alwan and H. V. Roberts, Time-series modeling for statistical process control, \emph{J. Bus. Econ. Statist.} \textbf{6} (1988), 87--95.
\bibitem{chatterjee} S. Chatterjee and P. Qiu, Distribution-free cumulative sum control charts using bootstrap-based control limits, \emph{Ann. Appl. Statist.} \textbf{3} (2009), 349--369.
\bibitem{edgington} E. S. Edgington and P. Onghena, \emph{Randomization Tests}, 4th ed., Chapman and Hall/CRC, 2007.
\bibitem{efron} B. Efron and R. J. Tibshirani, \emph{An Introduction to the Bootstrap}, Chapman and Hall, 1993.
\bibitem{guyatt} G. H. Guyatt et al., Determining optimal therapy: randomized trials in individual patients, \emph{New England J. Med.} \textbf{314} (1986), 889--892.
\bibitem{kunsch} H. R. K\"unsch, The jackknife and the bootstrap for general stationary observations, \emph{Ann. Statist.} \textbf{17} (1989), 1217--1241.
\bibitem{lillie} E. O. Lillie et al., The n-of-1 clinical trial: the ultimate strategy for individualizing medicine?, \emph{Personalized Medicine} \textbf{8} (2011), 161--173.
\bibitem{lorden} G. Lorden, Procedures for reacting to a change in distribution, \emph{Ann. Math. Statist.} \textbf{42} (1971), 1897--1908.
\bibitem{montgomery} D. C. Montgomery, \emph{Introduction to Statistical Quality Control}, 7th ed., Wiley, 2013.
\bibitem{page} E. S. Page, Continuous inspection schemes, \emph{Biometrika} \textbf{41} (1954), 100--115.
\bibitem{stone} M. Stone, Cross-validatory choice and assessment of statistical predictions, \emph{J. Roy. Statist. Soc. B} \textbf{36} (1974), 111--147.
\bibitem{wilson} E. B. Wilson, Probable inference, the law of succession, and statistical inference, \emph{J. Amer. Statist. Assoc.} \textbf{22} (1927), 209--212.
\end{thebibliography}

\end{document}
